\section{Introduction}

Language-model systems can answer from parameters, retrieve external evidence, cache reusable facts, or update themselves temporarily or durably. Existing work usually studies retrieval selection \citep{asai2024selfrag,yan2024crag}, memory use, or editing mechanisms \citep{wang2024wise,yu2023melo,zhang2024keStudy} in isolation. We instead study the controller problem: when should one policy choose among these operations under a single objective?

The preprint version of this project established two narrow claims. On PopQA \citep{mallen2022trust}, memory-aware routing can preserve recurring-case quality while reducing repeated retrieval. On a bundled deterministic FreshQA-style benchmark inspired by FreshLLMs/FreshQA \citep{vu2024freshllms}, the full six-action controller can decide when retrieval, memory, temporary adaptation, and guarded durable consolidation should be used for updated knowledge. Those claims remain intact and frozen.

This journal-oriented extension asks a harder question: do the same qualitative conclusions hold under broader evaluation? We therefore add seeded multi-model experiments, a broader PopQA run, a public FreshQA-derived benchmark built from weekly public releases, a public benchmark provenance package, a manual public-track audit, and a calibrated router targeted at the public changing-answer regime. The resulting picture is more valuable than a uniformly positive story. The recurring-memory result generalizes cleanly. The bundled six-action freshness result also remains strong. But the public FreshQA-derived benchmark is retrieval-dominant, and the main failure mode is not memory itself; it is delayed utility that is too aggressive for this regime.

That mixed result is precisely why the extension is useful. It upgrades the paper from a frozen artifact package to a broader empirical study with positive evidence, negative evidence, and a concrete calibrated fix. The paper is strongest not as a universal-win controller story, but as an evaluation-and-boundary study showing where multi-timescale routing works, where it fails, and how regime-aware delayed-utility gating materially improves the hardest public regime.
